
\chapter{Einleitung}

Unbemannte Fluggeräte werden zunehmend in zivilen Anwendungsbereichen eingesetzt.
Eine grundlegende Voraussetzung für ihren sicheren Betrieb ist die zuverlässige Kenntnis der eigenen Position.
Insbesondere bei einer autonomen Flugführung muss ein Fluggerät seine Lage im Raum fortlaufend bestimmen, um Kollisionen mit Hindernissen zu vermeiden und einen sicheren Flug zu gewährleisten.

Im Außenbereich wird die dafür benötigte Positionsreferenz in der Regel durch Satellitennavigation bereitgestellt.
In Innenräumen steht diese jedoch meist nicht zur Verfügung, da der Satellitenempfang durch die Gebäudestruktur eingeschränkt ist.
Ohne eine externe Positionsreferenz lässt sich die Position nur kurzzeitig aus den im Fluggerät verbauten Sensoren schätzen.
Dabei summieren sich Fehler mit der Zeit auf und führen im Innenraum zwangsläufig zu einer Kollision mit einem Hindernis.
Für einen zuverlässigen Flug im Innenraum ist daher eine alternative Positionsreferenz zwingend erforderlich.

Das Ziel dieser Arbeit ist es, ein unbemanntes Fluggerät zu entwickeln, das sich ohne Satellitennavigation im Innenraum selbst lokalisiert und positioniert.
Die extern berechnete Positionsschätzung soll dem Flugcontroller als Referenz für die Flugsteuerung bereitgestellt werden.
Darüber hinaus ist ein Verfahren zu entwickeln, das eine unbekannte Umgebung vollständig erfasst oder eine bestehende Karte ergänzt und deren Ergebnis der Positionsbestimmung zuführt.


Als Grundgerüst dient ein Multikopter des Typs Holybro X650 mit einem Flugcontroller Pixhawk 6X und einem als Begleitrechner eingesetzten NVIDIA Jetson Orin Nano.
Das Fluggerät ist mit zwei Logitech-Kameras ausgestattet, von denen eine nach unten und eine nach vorne gerichtet ist.
Die entwickelte Lösung ist vollständig aus quelloffener Software und gängigen Komponenten aufgebaut.

Als Grundlage der Positionsreferenz dienen Fiducial-Marker, die in der Flugumgebung angebracht sind.
Unter den betrachteten Fiducial-Marker-Systemen fiel die Wahl auf AprilTags des Typs \texttt{tag36h11}, da diese sich mit der eingesetzten Kartierungsmethode kombinieren lassen und hardwarebeschleunigt erkannt werden können.
Die Lage der Fiducial-Marker zueinander wird in einer vorab erstellten Karte hinterlegt.
Die Auswahl des Kartierungsverfahrens erfolgt über eine Nutzwertanalyse der betrachteten Verfahren zur simultanen Lokalisierung und Kartierung.


Während des Flugs berechnet ein auf dem Fluggerät befestigter Begleitrechner aus den Fiducial-Marker-Detektionen einer nach unten gerichteten Kamera fortlaufend die Pose des Fluggeräts relativ zur Karte.
Die einzelnen Schätzungen werden über eine Ausreißerfilterung bereinigt, anhand der Fiducial-Markergröße, des Abstands und der Güte der Karte gewichtet fusioniert und zeitlich geglättet.
Die so bestimmte Pose wird über eine Kommunikationsschnittstelle an den Flugcontroller übergeben.
Dort wird sie in einem zeitdiskreten Zustandsschätzer mit den übrigen Sensordaten fusioniert und als externe Positionsreferenz verwendet, während Satellitennavigation und Magnetometer von der Fusion ausgeschlossen sind.

Die Arbeit ordnet zunächst bestehende Verfahren zur Lokalisierung ohne Satellitennavigation in den Stand der Technik ein und leitet daraus die gewählte Marker- und Kartierungsmethodik ab.
Darauf folgen die Beschreibung des Gesamtsystems, die Umsetzung der Kartierung und die im Flug ablaufende Positionsbestimmung.
Den Abschluss bilden die Versuchsdurchführung und die Evaluation, in der die erzeugten Karten gegen eine bekannte Referenz verglichen und das Verhalten der Positionsbestimmung in praktischen Flugversuchen bewertet werden.